\documentclass[12pt,a4paper]{article}

\usepackage[french]{babel}
\usepackage[utf8]{inputenc}
\usepackage[T1]{fontenc}
\usepackage{fouriernc}
\usepackage{geometry}
\usepackage{array}
\usepackage[table]{xcolor}
\usepackage{amsmath}
\usepackage{fancyhdr}
\usepackage{tikz}
\usepackage{pgfplots}
\usepackage[most]{tcolorbox}
\usepackage{enumitem}
\pgfplotsset{compat=1.18}

\geometry{top=2.2cm, bottom=2cm, left=2cm, right=2cm, headheight=15pt}

%------------------------------------------------------------
% Couleurs
%------------------------------------------------------------
\definecolor{bluedark}{RGB}{50, 100, 160}
\definecolor{bluelight}{RGB}{210, 228, 248}
\definecolor{bluemid}{RGB}{170, 200, 235}
\definecolor{grisfond}{RGB}{245, 245, 245}
\definecolor{vertcorr}{RGB}{34, 120, 60}
\definecolor{vertlight}{RGB}{220, 245, 225}
\definecolor{vertmid}{RGB}{150, 210, 170}

%------------------------------------------------------------
% En-tête / pied de page
%------------------------------------------------------------
\pagestyle{fancy}
\fancyhf{}
\lhead{\small Physique-Chimie -- Terminale Bac Pro -- TRPM}
\rhead{\small \textbf{CORRIGÉ} -- Propagation du son}
\cfoot{\thepage}
\renewcommand{\headrulewidth}{0.5pt}

%------------------------------------------------------------
% Boîtes tcolorbox
%------------------------------------------------------------
\tcbset{
  stylepartie/.style={
    colback=bluedark, colframe=bluedark, coltext=white,
    boxrule=0pt, arc=3pt,
    left=6pt, right=6pt, top=4pt, bottom=4pt,
    before skip=10pt, after skip=6pt,
  },
  stylecorr/.style={
    colback=vertlight, colframe=vertmid,
    boxrule=0.8pt, arc=4pt,
    left=8pt, right=8pt, top=6pt, bottom=6pt,
    before skip=4pt, after skip=6pt,
  },
  styleinfo/.style={
    colback=grisfond, colframe=bluemid,
    boxrule=0.8pt, arc=3pt,
    left=8pt, right=8pt, top=5pt, bottom=5pt,
    before skip=6pt, after skip=6pt,
  },
}

%------------------------------------------------------------
% Commandes
%------------------------------------------------------------
\newcounter{numquestion}
\newcommand{\question}[1]{%
  \stepcounter{numquestion}%
  \vspace{8pt}%
  \noindent{\color{bluedark}$\blacktriangleright$~\textbf{Question~\thenumquestion.}}~#1%
  \vspace{2pt}%
}

\newcommand{\repcorr}[1]{%
  \begin{tcolorbox}[stylecorr]
    {\color{vertcorr} #1}
  \end{tcolorbox}%
}

%------------------------------------------------------------
\begin{document}
%------------------------------------------------------------

%============================================================
% TITRE
%============================================================
\begin{tcolorbox}[
  colback=white, colframe=bluedark, boxrule=1.5pt, arc=4pt,
  before skip=0pt, after skip=8pt,
  left=10pt, right=10pt, top=8pt, bottom=8pt,
]
  \begin{center}
    {\color{bluedark}\Large\bfseries Propagation du son}\\[4pt]
    {\large Activité documentaire -- \textbf{\color{vertcorr}CORRIGÉ}}\\[6pt]
    \hrule height 0.4pt \vspace{6pt}
    \textbf{Durée :} 45 min \hfill \textbf{Classe :} Terminale Bac Pro TRPM \hfill \textbf{Total : 20 points}
  \end{center}
\end{tcolorbox}

\vspace{6pt}

%============================================================
% PARTIE A
%============================================================
\begin{tcolorbox}[stylepartie]
  Partie A -- Nature et propriétés du son \hfill {\normalfont\small (15 min)}
\end{tcolorbox}

%------------------------------------------------------------
\question{Le son peut-il se propager dans le vide ? Justifier en citant le document~1.}

\repcorr{
  \textbf{Non}, le son ne peut pas se propager dans le vide. Le document~1 précise :
  \textit{« Le son ne se propage pas dans le vide. »} C'est une onde \textbf{mécanique}
  qui nécessite un milieu matériel (solide, liquide ou gaz) pour se propager.
}

%------------------------------------------------------------
\question{Comparer la célérité du son dans les gaz, les liquides et les solides. Dans quel milieu le son se propage-t-il le plus vite ?}

\repcorr{
  D'après le document~1 :
  \begin{itemize}[noitemsep]
    \item Dans les \textbf{gaz} : célérité faible $\approx 340$~m/s (air à 20°C).
    \item Dans les \textbf{liquides} : célérité plus élevée $\approx 1\,480$~m/s (eau).
    \item Dans les \textbf{solides} : célérité encore plus élevée, jusqu'à $6\,300$~m/s (aluminium).
  \end{itemize}
  $\Rightarrow$ Le son se propage le plus vite dans les \textbf{solides}.
}

%------------------------------------------------------------
\question{La célérité du son dans l'air à 20°C est $v = 340$~m/s. Convertir cette valeur en km/h.}

\repcorr{
  \[
    v = 340~\text{m/s} \times 3{,}6 = \boxed{1\,224~\text{km/h}}
  \]
}

\vspace{4pt}

%============================================================
% PARTIE B
%============================================================
\begin{tcolorbox}[stylepartie]
  Partie B -- Célérité du son dans les matériaux et application industrielle \hfill {\normalfont\small (15 min)}
\end{tcolorbox}

%------------------------------------------------------------
\question{Relever dans le tableau la célérité du son dans l'acier et dans l'aluminium. Laquelle est la plus élevée ?}

\repcorr{
  D'après le tableau du document~2 :
  \begin{itemize}[noitemsep]
    \item Acier : $v = 5\,900$~m/s
    \item Aluminium : $v = 6\,300$~m/s
  \end{itemize}
  $\Rightarrow$ La célérité est plus élevée dans l'\textbf{aluminium} ($6\,300 > 5\,900$~m/s).
}

%------------------------------------------------------------
\question{Pourquoi parle-t-on de contrôle \textit{non destructif} ? Quel est l'intérêt pour un technicien TRPM ?}

\repcorr{
  Le contrôle est dit \textbf{non destructif} car on vérifie l'intégrité de la pièce
  \textbf{sans l'endommager} ni la détruire : les ondes ultrasonores traversent la pièce
  sans l'abîmer.\\[4pt]
  \textbf{Intérêt pour le technicien TRPM :} on peut contrôler la totalité de la
  production, et les pièces conformes peuvent être directement utilisées. Cela évite
  le gaspillage et garantit la sécurité des pièces en service.
}

%------------------------------------------------------------
\question{Dans la formule $d = \dfrac{v \times \Delta t}{2}$, expliquer pourquoi on divise par 2.}

\repcorr{
  L'onde ultrasonore effectue un trajet \textbf{aller-retour} : elle part du capteur,
  atteint l'interface (fond ou défaut), puis revient au capteur. La durée $\Delta t$
  correspond donc à deux fois la distance réelle. On divise par 2 pour obtenir la
  distance \textbf{simple} (capteur → interface).
}

\vspace{4pt}

%============================================================
% PARTIE C
%============================================================
\begin{tcolorbox}[stylepartie]
  Partie C -- Application : contrôle d'une pièce en acier \hfill {\normalfont\small (15 min)}
\end{tcolorbox}

%------------------------------------------------------------
\question{En utilisant l'écho B et la formule $d = \dfrac{v \times \Delta t_B}{2}$, calculer l'épaisseur de la pièce mesurée. Comparer au résultat attendu ($e = 60$~mm).}

\repcorr{
  \textbf{Données :} $v = 5\,900$~m/s \quad $\Delta t_B = 20{,}3~\mu\text{s} = 20{,}3 \times 10^{-6}$~s
  \[
    d = \frac{v \times \Delta t_B}{2} = \frac{5\,900 \times 20{,}3 \times 10^{-6}}{2}
      = \frac{0{,}11977}{2} \approx \boxed{0{,}0599~\text{m} = 59{,}9~\text{mm}}
  \]
  \textbf{Comparaison :} $59{,}9$~mm $\approx 60$~mm $\Rightarrow$ le résultat est cohérent avec
  l'épaisseur nominale (écart $< 0{,}2$~\%).
}

%------------------------------------------------------------
\question{En utilisant l'écho A, calculer la profondeur à laquelle se trouve l'anomalie interne dans la pièce.}

\repcorr{
  \textbf{Données :} $v = 5\,900$~m/s \quad $\Delta t_A = 8{,}1~\mu\text{s} = 8{,}1 \times 10^{-6}$~s
  \[
    d = \frac{v \times \Delta t_A}{2} = \frac{5\,900 \times 8{,}1 \times 10^{-6}}{2}
      = \frac{0{,}04779}{2} \approx \boxed{0{,}0239~\text{m} = 23{,}9~\text{mm}}
  \]
  L'anomalie se situe à environ \textbf{23,9~mm} de la surface de la pièce.
}

%------------------------------------------------------------
\question{\textbf{Synthèse :} Le technicien décide de rejeter cette pièce. Justifier cette décision à partir des résultats obtenus aux questions précédentes.}

\repcorr{
  L'écho A révèle la présence d'une \textbf{anomalie interne} (fissure, inclusion ou
  porosité) à 23,9~mm de profondeur dans la pièce. Une telle anomalie peut
  \textbf{fragiliser la pièce} et provoquer une rupture prématurée en service,
  ce qui représente un risque de sécurité inacceptable.\\[4pt]
  $\Rightarrow$ Le technicien a raison de \textbf{rejeter la pièce} : elle n'est pas conforme
  au cahier des charges (absence de défaut interne).
}

\vspace{10pt}
\begin{center}
  \color{bluedark}\rule{0.5\linewidth}{0.6pt}
\end{center}

\end{document}
